\begin{table}[H]
\centering
\begin{threeparttable}
\caption{Illustrative Quiz Battery and Standardized Scoring for CTRL\_021}
\label{tab:illustrative_quiz}
\scriptsize
\setlength{\tabcolsep}{4pt}
\begin{tabularx}{\textwidth}{p{2.6cm}X p{3.8cm} p{1.3cm}}
\toprule
Quiz type & Illustrative prompt & Expected answer / scoring focus & Points \\
\midrule
Reading comprehension (\textit{yuedu lijie}) & Why does the article treat the Louvre theft as more than an ordinary burglary? & Full credit requires identifying both the symbolic stakes and the institutional stakes: the theft targeted crown jewels in France's most visible museum and exposed a major security failure. & 20 \\
\addlinespace
Grammar / syntax (\textit{yufa jufa}) & In the sentence ``That a truck could park outside the museum without raising suspicion is worrying enough,'' what grammatical role does the initial \textit{that}-clause play? & Subject clause. Full credit requires recognizing that the clause functions as the grammatical subject of the sentence rather than as reported speech or an object clause. & 20 \\
\addlinespace
Vocabulary (\textit{danci ceshi}) & In context, what does \textit{make off with} most nearly mean? & ``Escape with stolen items'' or an equivalent paraphrase. & 20 \\
\addlinespace
Listening (\textit{tingli ceshi}) & After the short instructor recap, what is the main lesson for museum security? & Full credit requires identifying the delay logic emphasized in the recap: effective security must slow thieves down and buy time for response, not rely on alarms alone. & 20 \\
\addlinespace
Sentence accumulation (\textit{juzi jilei}) & Complete the model sentence used in the lesson review: ``The point of museum security is not only to protect valuable objects but also to \_\_\_\_\_ long enough for guards to respond.'' & ``Buy time'' or a grammatically correct equivalent such as ``slow thieves down.'' Full credit requires preserving the sentence's causal logic and usable phrase structure. & 20 \\
\bottomrule
\end{tabularx}
\begin{tablenotes}[flushleft]
\footnotesize
\item Note: This table provides one illustrative item for each observed quiz subtype in the app. In the teaching data, subtype scores were standardized to run from 0 to 20 and then summed to a 0--100 total for the lesson-level quiz score. The prompts shown here are paraphrased from the adapted lesson and video recap rather than reproduced from copyrighted source text.
\end{tablenotes}
\end{threeparttable}
\end{table}
